% =============================================================================
% Chapter 1 --- Introduction
% =============================================================================

\chapter{Introduction}
\label{ch:introduction}

\section{Motivation: The Memory Wall and In-Memory Computing}
\label{sec:intro-motivation}

The energy and latency costs of moving data between memory and arithmetic
units dominate modern computing workloads, ANN inference and training
foremost among them. This \emph{memory wall} is intrinsic to any
architecture that keeps storage and computation
separate~\cite{Ielmini2018RRAM, Verma2019IMC, Sebastian2020IMC}, and the
penalty is largest for ANN workloads whose dominant kernel ---
matrix--vector multiplication (MVM) over large weight matrices --- is
data-movement-bound.

\emph{In-memory computing} (IMC) confronts the bottleneck by performing
the computation directly inside the memory array, exploiting both the
physical properties of the device and the parallelism of array-level
organization~\cite{Sebastian2020IMC, Verma2019IMC}. Non-volatile
resistive-switching devices --- RRAM, PCM, and floating-gate memristors
--- can store a quantized synaptic weight as a programmable conductance,
turning a 2D crossbar into a one-shot MVM
engine~\cite{Ielmini2018RRAM, Wang2022YFlash}.

\section{Project Context and Scope}
\label{sec:intro-context}

This report documents the Technion VLSI Laboratory final project
\textbf{``Design Controller for an Artificial Neural Network (ANN)
Accelerator and its Peripherals''}. The project addresses the digital
control layer that surrounds an analog IMC core: a controller that
decodes host traffic, sequences program, read, erase, and inference
operations, and drives an analog ANN core based on Y-Flash floating-gate
memristors through a packed transaction word and a multi-bit pulse-mode
bus. The target inference workload for the accelerator is image
classification on the MNIST dataset --- a standard benchmark of
$28\times28$ grayscale images of handwritten digits ($0$--$9$), with
each pixel quantized to one byte. The $784$ input pixels of an MNIST
image are streamed eight at a time into the input buffer
(\Cref{sec:rtl-buffer}) and consumed by the analog core through eight
bit-serial lanes, so the choice of MNIST sets the size of the
per-cycle activation window and the bit-width of the controller's
collect/compute datapath.

The design is motivated by a direct consequence of the device physics: a
single \emph{write-once} operation cannot be relied upon, because
device-to-device and cycle-to-cycle variability displace the resulting
conductance from the target. A closed-loop
\emph{program--verify--(re-program or erase)} control policy is therefore
required to drive stored weights toward their intended
values~\cite{Ielmini2018RRAM, Sebastian2020IMC, Wang2022YFlash}. The
central piece of digital control is consequently a hierarchical FSM that
mediates between the host and the analog core and that implements this
recovery logic. The project scope is confined to RTL implementation and
functional verification; netlist generation, timing closure, and
power/area analysis are out of scope.

\section{Contributions}
\label{sec:intro-contributions}

The engineering contributions of this work are organized along three
axes:

\begin{itemize}
    \item \textbf{A controller-centric RTL architecture}, factored into a
          combinational host interface, a hierarchical FSM controller with
          three cooperating sub-FSMs for programming, verification, and
          erase, and a synchronous input buffer that stages inference
          pixels and stores expected weights for the verify path. The
          controller implements a data-driven pulse-shaping mechanism in
          which the first programming attempt for each cell is derived
          from a look-up table (LUT) indexed by the expected weight.
    \item \textbf{A closed-loop recovery policy} that resolves the
          write--verify problem at the digital control layer. After every
          programming attempt the controller reads the cell back,
          compares the readback against the expected weight, and branches
          into a short retry pulse train (under-programmed), an erase
          followed by reprogramming (over-programmed, or after retry
          exhaustion), or termination on a successful verify.
    \item \textbf{A directed verification environment} that exercises
          this recovery policy without access to the physical analog
          core. The environment models the memristor array with a
          behavioral shadow weight matrix, generates an analog-side
          handshake that follows the controller's own sub-FSM state, and
          deliberately injects under- and over-programmed readbacks
          during the verify window so that the recovery branches are
          observed in simulation. This is the methodological idea that
          lets the project deliver high confidence in the RTL prior to
          integration with the real analog ANN.
\end{itemize}

\section{Report Organization}
\label{sec:intro-organization}

\Cref{ch:background} reviews IMC, the Y-Flash device, and the
write--verify problem that motivates the controller architecture.
\Cref{ch:architecture} presents the three-module stack, the transaction
word, and the host command set. \Cref{ch:rtl} is the technical heart of
the report: each module is described as an architectural unit, with
emphasis on the main FSM, the three sub-FSMs, and the LUT-based pulse
shaper. \Cref{ch:verification} documents the verification strategy and
the behavioral mock plus fault-injection methodology.
\Cref{ch:results} consolidates the simulation evidence.
\Cref{ch:limitations} catalogs the limitations of the present work and
proposes a roadmap for follow-on verification. \Cref{ch:conclusion}
closes the report. Comprehensive supporting material --- full FSM
diagrams, module port tables, addressing/packing conventions, pulse
timing arithmetic, testbench mechanics, and RTL source listings --- is
collected in the project repository (see Supplementary Material at the
end of the report).
